\documentclass[11pt]{article}
\usepackage{tcee}
\addbibresource{3B42.bib}
\tcetema[Relaciones exteriores y cooperación al desarrollo de la UE]{3.B.42}{Las relaciones económicas exteriores de la Unión Europea. La política de cooperación al desarrollo de la Unión Europea.}
\tcefecha{14/11/2026}
\begin{document}
\tceetitulo

\section*{Introducción}
\minutos{2,5}

\begin{itemize}
\item \textbf{\emph{Enganche}}:
  \begin{itemize}
  \item «Más vale vecino cercano que hermano lejano», dice el refrán. En 2025 la Unión Europea lo comprobó con su aliado de siempre: Estados Unidos le impuso un arancel del 15~\% a la mayoría de sus exportaciones y recortó su ayuda al desarrollo a menos de la mitad\nocite{UE2025DeclaracionEEUU,OCDE2026AODPreliminar}.
  \item Ese mismo año la Unión siguió siendo el primer donante del mundo, presentó un Pacto por el Mediterráneo con sus vecinos del Sur y se preparó para abrir las negociaciones de adhesión con Ucrania y Moldavia, que empezaron en 2026\nocite{CE2025PactoMediterraneo,CE2026PaqueteAmpliacion}.
  \item La pregunta que recorre la exposición es: \emph{¿para qué usa la Unión su peso económico fuera de sus fronteras, y le basta con comerciar para ser un actor en el mundo?}
  \end{itemize}

\item \textbf{\emph{Relevancia}}:
  \begin{itemize}
  \item \emph{Es una potencia económica}: la Unión es la segunda economía del mundo a tipos de cambio de mercado, el primer exportador e importador de servicios, el segundo de bienes y el primer socio comercial de unos 80 países.
  \item \emph{Es la primera potencia de la ayuda}: el \concepto{Equipo Europa} (\emph{Team Europe}), la Unión con sus Estados miembros, da la mitad de la ayuda oficial al desarrollo de los países donantes de la OCDE.
  \item \emph{Para España}: sus relaciones con el Mediterráneo y con América Latina pasan en buena parte por Bruselas, y su propia cooperación se coordina con la de la Unión \vertema{4.A.26}.
  \end{itemize}

\item \textbf{\emph{Contextualización}}:
  \begin{itemize}
  \item \emph{Qué son}: las \concepto{relaciones económicas exteriores} de la Unión son el conjunto de políticas con las que se relaciona con terceros países y organizaciones. No hay una sola política, sino varias: la comercial, la de ampliación, la de vecindad, la cooperación al desarrollo, la ayuda humanitaria, la ayuda macrofinanciera y las sanciones, más la representación exterior del euro.
  \item \emph{Cómo se reparten las competencias}: la política comercial es exclusiva de la Unión; la cooperación al desarrollo y la ayuda humanitaria son competencias \emph{paralelas}, que ejercen a la vez la Unión y los Estados; la política exterior y de seguridad común (PESC) se decide casi siempre por unanimidad \vertema{3.B.36}.
  \item \emph{El objeto de esta exposición}: el marco, la geografía y los socios de esas relaciones y, con detalle, la cooperación al desarrollo. La política comercial tiene su propio tema \vertema{3.B.41}, y la financiación del desarrollo en general, también \vertema{3.B.29}.
  \end{itemize}

\item \textbf{\emph{Problemática}}:
  \begin{itemize}
  \item ¿Con qué reglas, actores e instrumentos actúa la Unión fuera de sus fronteras, y cuánto pesa en el mundo?
  \item ¿Por qué trata de forma distinta a sus candidatos, a sus vecinos y a sus grandes socios, y le funciona?
  \item ¿Qué hace la Unión por el desarrollo, con qué dinero y con qué resultados?
  \end{itemize}

\item \textbf{\emph{Estructura}}:
  \begin{itemize}
  \item La exposición va de lo general a lo particular, y del mundo a los países más pobres:
    \begin{itemize}
    \item \concepto{I. Visión global} (\ref{sec:global}): el marco jurídico e institucional, el peso económico de la Unión en el mundo, con el modelo del efecto Bruselas, y su lugar en las organizaciones internacionales.
    \item \concepto{II. Visión regional} (\ref{sec:regional}): la ampliación, con el modelo de la condicionalidad; la vecindad; las sanciones, con el modelo del tope al precio del petróleo ruso; y los grandes socios.
    \item \concepto{III. La cooperación al desarrollo} (\ref{sec:cooperacion}): conceptos y justificación, con el modelo del reparto eficiente de la ayuda, y la política de la Unión: datos, marco, historia, instrumentos y retos.
    \end{itemize}
  \end{itemize}
\end{itemize}

\begin{esquema}[unbreakable]
\textbf{3.B.42. Relaciones económicas exteriores y cooperación al desarrollo de la UE}\\[2pt]
Introducción\esqtiempo{2,5}\\
I. Visión global: marco y peso de la Unión\esqtiempo{7,5}\\
\hspace*{1.5em}1. El marco jurídico e institucional\\
\hspace*{1.5em}2. El peso económico de la Unión en el mundo\\
\hspace*{1.5em}3. Las organizaciones y foros internacionales\\
II. Visión regional: vecinos, socios y rivales\esqtiempo{11}\\
\hspace*{1.5em}1. La economía como instrumento geoestratégico\\
\hspace*{1.5em}2. El movimiento de bienes y factores: los grandes socios\\
III. La política de cooperación al desarrollo\esqtiempo{7,5}\\
\hspace*{1.5em}1. Consideraciones generales\\
\hspace*{1.5em}2. La Unión Europea y la ayuda al desarrollo\\
Conclusión\esqtiempo{1,5}
\end{esquema}

\begin{notaopositor}
Reparto del cante: introducción, \tiempo{2,5}; I, \tiempo{7,5} (I.1, \tiempo{2,5}; I.2, \tiempo{3,5}, de ellos \tiempo{2} para el efecto Bruselas; I.3, \tiempo{1,5}); II, \tiempo{11} (II.1, \tiempo{7,5}: ampliación con la condicionalidad, \tiempo{3,5}; vecindad, \tiempo{1,5}; sanciones con el tope al petróleo, \tiempo{2,5}; II.2, \tiempo{3,5}); III, \tiempo{7,5} (III.1, \tiempo{2,5}, con el reparto eficiente de la ayuda; III.2, \tiempo{5}), y conclusión, \tiempo{1,5}. Los gráficos y figuras imprescindibles son los círculos de las relaciones exteriores, el de la condicionalidad, el del tope al petróleo y el del reparto de la ayuda, y una cifra de la ayuda del Equipo Europa. Lo que va en cajas grises, en notas y en los anexos es para las preguntas del tribunal.
\end{notaopositor}

\begin{anotaciones}
\textbf{Notación y siglas.} En el efecto Bruselas, $M_{UE}$ y $M_R$ son las ventas de una empresa en la Unión y en el resto del mundo, $c$ el coste por unidad de cumplir la norma europea y $F$ el coste fijo de mantener 2 líneas de producción. En la condicionalidad, $B$ es el valor de la adhesión para el candidato, $p$ la probabilidad de que llegue (credibilidad), $T$ los años que faltan, $\delta$ el factor de descuento, $C$ el coste interno de las reformas y $g$ las recompensas intermedias. En el tope al petróleo, $P^*$ es el precio mundial, $\bar P$ el tope, $c$ el coste marginal, $\bar Q$ la capacidad de exportación y $\theta$ la parte del crudo que pasa por servicios occidentales. En el reparto de la ayuda, $A_i$ es la ayuda al país $i$ y $m_i$ la reducción marginal de la pobreza. TUE y TFUE: Tratado de la Unión Europea y Tratado de Funcionamiento de la Unión Europea; AOD: ayuda oficial al desarrollo; CAD: Comité de Ayuda al Desarrollo de la OCDE; RNB: renta nacional bruta; AR: Alta Representante de la Unión para Asuntos Exteriores y Política de Seguridad; IVDCI: Instrumento de Vecindad, Desarrollo y Cooperación Internacional.
\end{anotaciones}

% ==========================================================================
\section{Visión global: marco y peso de la Unión}\label{sec:global}
\minutos{7,5}

\begin{itemize}
\item Empezamos por el marco común a todas las relaciones exteriores: quién decide, con qué reglas y con qué peso.
\item \emph{Conclusión anticipada}: la Unión actúa fuera con muchas manos (la Comisión, el Consejo, la Alta Representante y 27 Estados) y casi solo con instrumentos económicos: el acceso a su mercado, sus normas, su dinero y su moneda. Ese es su poder, y también su límite.
\end{itemize}

% --------------------------------------------------------------------------
\subsection{El marco jurídico e institucional}\label{sec:marco}

\subsubsection{La base jurídica}\label{sec:base}

\begin{itemize}
\item \emph{Los principios y objetivos} (arts. 21 y 22 TUE): la acción exterior se guía por los principios que inspiraron la creación de la Unión (democracia, Estado de Derecho, derechos humanos, igualdad, solidaridad, respeto del Derecho internacional) y persigue, entre otros objetivos\nocite{UE2016TUE}:
  \begin{itemize}
  \item apoyar el desarrollo sostenible de los países en desarrollo, con el objetivo fundamental de erradicar la pobreza;
  \item fomentar la integración de todos los países en la economía mundial, también mediante la supresión progresiva de los obstáculos al comercio;
  \item ayudar a las poblaciones afectadas por catástrofes y promover un sistema internacional basado en una cooperación multilateral sólida y en una buena gobernanza mundial.
  \end{itemize}
  El Consejo Europeo fija por unanimidad los intereses y objetivos estratégicos de la Unión (art. 22 TUE).
\item \emph{La quinta parte del TFUE} (arts. 205 a 222) reúne toda la acción exterior que no es PESC (tabla~\ref{tab:tfue})\nocite{UE2016TFUE}.
\item \emph{Otras bases}: la política de vecindad (art. 8 TUE), la adhesión (art. 49 TUE) y la PESC (arts. 23 a 46 TUE).
\end{itemize}

\begin{table}[H]
\centering
\caption{La acción exterior en el TFUE (quinta parte, arts. 205-222)}\label{tab:tfue}
\small
\begin{tabular}{>{\raggedright\arraybackslash}m{2.4cm}>{\raggedright\arraybackslash}m{5.0cm}>{\raggedright\arraybackslash}m{6.6cm}}
\toprule
\rowcolor{tceeazulpalido} Artículos & Materia & Qué dice \\
\midrule
205 & Disposiciones generales & La acción exterior sigue los principios y objetivos del art. 21 TUE \\
206-207 & Política comercial común & Competencia exclusiva; aranceles, acuerdos comerciales, defensa comercial \vertema{3.B.41} \\
208-211 & Cooperación para el desarrollo & Objetivo principal: reducir y erradicar la pobreza; políticas de la Unión y de los Estados que se complementan \\
212-213 & Cooperación económica, financiera y técnica con terceros países & Con los países que no son en desarrollo: base de la ayuda macrofinanciera; ayuda financiera urgente (art. 213) \\
214 & Ayuda humanitaria & Asistencia a las víctimas de catástrofes y conflictos; Cuerpo Voluntario Europeo de Ayuda Humanitaria \\
215 & Medidas restrictivas & Sanciones económicas, tras una decisión PESC por unanimidad \\
216-219 & Acuerdos internacionales & Quién negocia y concluye; acuerdos de asociación (art. 217); acuerdos cambiarios (art. 219) \\
220-221 & Organizaciones internacionales y delegaciones & Cooperación con las Naciones Unidas, la OCDE y otras; delegaciones de la Unión \\
222 & Cláusula de solidaridad & Ayuda mutua ante un ataque terrorista o una catástrofe \\
\bottomrule
\end{tabular}
\fuente{elaboración propia a partir del TFUE\nocite{UE2016TFUE}.}
\end{table}

\subsubsection{Las competencias y los actores}\label{sec:actores}

\begin{itemize}
\item \emph{Las competencias son de 3 tipos} (anexo~\ref{anexo:competencias}):
  \begin{itemize}
  \item \emph{exclusivas}: la política comercial y la unión aduanera (art. 3 TFUE); la Unión habla sola y los Estados no pueden firmar acuerdos comerciales propios;
  \item \emph{paralelas}: en la cooperación al desarrollo y la ayuda humanitaria, la Unión tiene su propia política, «sin que el ejercicio de esta competencia pueda tener por efecto impedir a los Estados miembros ejercer la suya» (art. 4.4 TFUE); por eso conviven 28 políticas de desarrollo, que deben coordinarse (art. 210 TFUE);
  \item \emph{intergubernamentales}: la PESC, decidida por unanimidad, sin actos legislativos y casi sin control del Tribunal de Justicia.
  \end{itemize}
\item \emph{Los actores} (figura~\ref{figu:actores}):
  \begin{itemize}
  \item el \emph{Consejo Europeo} fija las orientaciones (presidido desde diciembre de 2024 por \autor{António Costa})\nocite{CE2026Organigrama};
  \item el \emph{Consejo de Asuntos Exteriores} decide los acuerdos, las sanciones y las posiciones; lo preside la AR, salvo cuando se reúne como Consejo de Comercio;
  \item la \emph{Alta Representante y vicepresidenta de la Comisión}, \autor{Kaja Kallas} desde diciembre de 2024, con el \concepto{Servicio Europeo de Acción Exterior} (SEAE) y sus 145 delegaciones\nocite{SEAE2026};
  \item la \emph{Comisión}, que propone, negocia los acuerdos comerciales y gestiona el dinero\nocite{CE2026Organigrama}: la DG de Comercio y Seguridad Económica (DG TRADE), la de Asociaciones Internacionales (DG INTPA, desarrollo), la de Ampliación y Vecindad Oriental (DG ENEST), la de Oriente Próximo, Norte de África y Golfo (DG MENA); las 2 últimas nacieron en febrero de 2025 al partirse la antigua DG de Vecindad y Ampliación, la de Protección Civil y Ayuda Humanitaria (DG ECHO) y la de Asuntos Económicos (DG ECFIN, ayuda macrofinanciera);
  \item el \emph{Parlamento Europeo}, que aprueba los acuerdos (art. 218 TFUE), es colegislador de los instrumentos financieros y autoridad presupuestaria;
  \item y los bancos públicos: el \emph{BEI Global} y el BERD \vertema{3.B.30}.
  \end{itemize}
\end{itemize}

\begin{infografia}[H]
\centering
\caption{Quién decide la acción exterior de la Unión}\label{figu:actores}
\grafico[0.95\textwidth]{actores-accion-exterior}
\fuente{elaboración propia a partir del TUE, el TFUE y el organigrama de la Comisión\nocite{UE2016TUE,UE2016TFUE,CE2026Organigrama}.}
\end{infografia}

\subsubsection{Los instrumentos}\label{sec:instrumentos}

\begin{itemize}
\item \emph{Los acuerdos internacionales} (art. 218 TFUE): la Comisión (o la AR, si es PESC) recomienda, el Consejo autoriza y fija las directrices, la Comisión negocia, el Parlamento da su aprobación y el Consejo concluye. Son de 4 grandes tipos:
  \begin{enumerate}
  \item \emph{de asociación} (art. 217 TFUE): diálogo político, cooperación y una zona de libre comercio (los de vecindad, América Latina, el Acuerdo de Samoa);
  \item \emph{comerciales}: zonas de libre comercio de nueva generación, con servicios, inversión y desarrollo sostenible \vertema{3.B.41};
  \item \emph{de asociación económica} con los países de África, el Caribe y el Pacífico (ACP), compatibles con la OMC;
  \item \emph{de colaboración y cooperación}, sin preferencias comerciales.
  \end{enumerate}
  Si tocan competencias de los Estados (la inversión de cartera, la solución de diferencias inversor-Estado), son \concepto{acuerdos mixtos} y los ratifican también los 27, como dejó claro el Dictamen 2/15 del Tribunal de Justicia (2017) sobre el acuerdo con Singapur\nocite{TJUE2017Dictamen215}.
\item \emph{Los instrumentos financieros}, en la rúbrica 6 del marco financiero plurianual, «Vecindad y mundo» (113.700 millones de euros en 2021-2027, el 9~\% del marco)\nocite{CE2026Rubricas} \vertema{3.B.37}:
  \begin{itemize}
  \item el IVDCI-Europa Global (79.462 millones), que reúne desde 2021 los antiguos instrumentos geográficos y temáticos y el Fondo Europeo de Desarrollo\nocite{UE2021IVDCI};
  \item el Instrumento de Ayuda Preadhesión, IPA III (14.162 millones)\nocite{UE2021IPA};
  \item la ayuda humanitaria (11.569 millones), el presupuesto de la PESC y los países y territorios de ultramar;
  \item y, fuera de los techos, el Mecanismo para Ucrania (50.000 millones en 2024-2027)\nocite{CE2026UcraniaDesembolso}.
  \end{itemize}
\item \emph{La ayuda macrofinanciera} (art. 212 TFUE): préstamos o donaciones a países vecinos con problemas de balanza de pagos, complementarios de un programa del FMI y condicionados a reformas.
\item \emph{Las medidas restrictivas} (art. 215 TFUE): congelación de activos, prohibiciones de importar o exportar y de prestar servicios, tras una decisión PESC unánime (\ref{sec:sanciones}).
\end{itemize}

% --------------------------------------------------------------------------
\subsection{El peso económico de la Unión en el mundo}\label{sec:peso}

\subsubsection{Comercio, inversión, ayuda y moneda}\label{sec:datos-peso}

\begin{itemize}
\item \emph{La producción}: la Unión tiene unos 450 millones de habitantes, el 5,5~\% de la población mundial, y en 2025 produjo el 17,9~\% del PIB mundial a tipos de cambio de mercado, por detrás de Estados Unidos (26,0~\%) y por delante de China (16,5~\%); en paridad de poder adquisitivo, el 14,0~\%\nocite{BancoMundial2026PIB}.
\item \emph{El comercio} (2025, sin el comercio entre sus Estados): es el segundo exportador de bienes del mundo (13,8~\%, tras China) y el segundo importador (12,9~\%, tras Estados Unidos), y el primer exportador e importador de servicios comerciales (en torno al 22,7~\% del total)\nocite{OMC2026Estadisticas}. Es el primer socio comercial de unos 80 países\nocite{CE2026PosicionComercio}.
  \begin{itemize}
  \item En bienes, sus 3 grandes socios son Estados Unidos, China y el Reino Unido (gráfico~\ref{fig:comercio-socios}): Estados Unidos es el primer cliente y China, el primer proveedor, con un superávit europeo de unos 197.000 millones de euros con el primero y un déficit de unos 375.000 millones con el segundo en 2025.
  \item En servicios, Estados Unidos, el Reino Unido y Suiza.
  \end{itemize}
\item \emph{La inversión}: a finales de 2024, las empresas de la Unión tenían 9,3 billones de euros invertidos fuera de ella, casi un 30~\% en Estados Unidos, y las de fuera, 7,0 billones en la Unión\nocite{Eurostat2025IED}.
\item \emph{La ayuda}: los Estados de la Unión dieron en 2025 la mitad de la AOD de los países del CAD, y las instituciones de la Unión, otros 23.100 millones de euros\nocite{OCDE2026DAC1} (\ref{sec:aod-ue}).
\item \emph{La moneda}: el euro es la segunda moneda internacional, a mucha distancia del dólar: el 20,2~\% de las reservas de divisas del mundo, frente al 57~\% del dólar, y un 60~\% de las exportaciones de la zona del euro fuera de ella se facturan en euros\nocite{BCE2026PapelEuro}. Su uso internacional da a la Unión señoreaje, menos riesgo cambiario para sus empresas y peso geopolítico, y le exige una política monetaria que tenga en cuenta al resto del mundo (anexo~\ref{anexo:euro}) \vertema{3.B.43}.
\end{itemize}

\begin{figure}[H]
\centering
\caption{Comercio de bienes de la Unión con sus principales socios, 2025 (miles de millones de euros)}\label{fig:comercio-socios}
\grafico[0.9\textwidth]{comercio-socios}
\fuente{elaboración propia a partir de Eurostat\nocite{Eurostat2026ComercioSITC}. Comercio extracomunitario de la UE-27.}
\end{figure}

\begin{itemize}
\item Pero el peso de la Unión no está solo en sus cifras. Fuera, se la conoce sobre todo por sus normas: ¿por qué acaban aplicándolas empresas y países que no forman parte de ella?
\end{itemize}

\subsubsection{El efecto Bruselas}\label{sec:bruselas}

\paragraph{a) Idea}

\begin{itemize}
\item La Unión no tiene un gran ejército ni la moneda dominante, pero fija normas que se aplican en todo el mundo: la protección de datos, la seguridad de los productos químicos, los alimentos o la competencia digital.
\item \textcite{Bradford2012Brussels} llamó \concepto{efecto Bruselas} a esa exportación unilateral de normas: la Unión no las negocia con nadie; son las empresas las que deciden aplicarlas en todas partes porque les sale más barato que tener 2 versiones de su producto\nocite{Bradford2020Brussels}.
\item Es la versión económica del \concepto{poder normativo} de la Unión (\emph{normative power Europe}), un poder de atracción y no de coerción, como el \concepto{poder blando} (\emph{soft power}) de \textcite{Nye2004Soft}\nocite{Manners2002Normative}.
\end{itemize}

\paragraph{b) Modelo}

\subparagraph{i) Supuestos}

\begin{enumerate}
\item \emph{Una empresa} vende $M_{UE}$ unidades en la Unión y $M_R$ en el resto del mundo, con un margen por unidad $\pi$ antes de cumplir normas.
\item \emph{La Unión impone una norma más estricta} que la del resto del mundo: cumplirla cuesta $c$ por unidad.
\item \emph{2 opciones de producción}: una sola línea con la norma europea para todo el mundo, o 2 líneas (la europea y la menos exigente), con un coste fijo adicional $F$ (diseño, procesos, cumplimiento y control separados).
\item \emph{La norma afecta a un objetivo poco móvil}: el consumidor europeo (no la empresa) es el que está en la Unión, y la empresa no puede venderle desde fuera sin cumplirla.
\end{enumerate}

\subparagraph{ii) Desarrollo}

\begin{itemize}
\item \emph{Paso 1: ¿se queda la empresa en el mercado europeo?}
  \begin{itemize}
  \item Sí, si lo que gana allí cubre el coste de cumplir: $(\pi-c)\,M_{UE}\geq0$. Con un mercado grande y rico, casi siempre es así.
  \end{itemize}
\item \emph{Paso 2: ¿una línea o 2?}
  \begin{itemize}
  \item Con una sola línea europea, paga $c$ también por lo que vende fuera: un coste $c\,M_R$.
  \item Con 2 líneas, se ahorra ese coste pero paga $F$.
  \end{itemize}
\item \emph{Paso 3: la regla}. La empresa adopta la norma europea en todo el mundo si
  \funcionclave[eq:bruselas]{\text{efecto Bruselas}\iff F> c\,M_R}
  es decir, si producir por separado cuesta más que cumplir la norma también fuera (gráfico~\ref{fig:bruselas}).
\end{itemize}

\begin{figure}[H]
\centering
\caption{El efecto Bruselas: cuándo extiende una empresa la norma de la UE a todo el mundo}\label{fig:bruselas}
\grafico[0.62\textwidth]{efecto-bruselas}
\fuente{elaboración propia a partir de \textcite{Bradford2012Brussels}.}
\end{figure}

\subparagraph{iii) Implicaciones}

\begin{itemize}
\item \emph{Las 5 condiciones de \autor{Anu Bradford}} salen del modelo:
  \begin{enumerate}
  \item un \emph{mercado grande} ($M_{UE}$ alto: nadie renuncia a él);
  \item \emph{capacidad reguladora} (instituciones que hacen y aplican normas complejas);
  \item \emph{normas estrictas} (si fueran laxas, no habría nada que exportar);
  \item \emph{objetivos poco móviles} (se regula al consumidor, no al capital, que se iría);
  \item \emph{indivisibilidad} de la producción ($F$ alto).
  \end{enumerate}
\item \emph{Del efecto de hecho al de derecho}: cuando las empresas ya cumplen la norma europea, les interesa que su país la adopte para no estar en desventaja frente a sus competidores locales; así se extendió el Reglamento General de Protección de Datos (RGPD) a otras legislaciones.
\item \emph{Lo que explica}: por qué la Unión influye más con sus normas que con su política exterior clásica.
\end{itemize}

\paragraph{c) Extensiones}

\begin{itemize}
\item \emph{La geoeconomía}: desde 2020 la Unión usa su mercado también como palanca defensiva, con instrumentos como el anticoerción (2023), el control de inversiones extranjeras o el mecanismo de ajuste en frontera por carbono, que mezclan la lógica del efecto Bruselas con la de la seguridad económica \vertema{3.B.41} \vertema{3.B.31}.
\end{itemize}

\paragraph{d) Evidencia empírica}

\begin{itemize}
\item \emph{Casos}: tras el RGPD (2018), grandes empresas tecnológicas extendieron sus protecciones a todos sus usuarios, y muchos países (Japón, Brasil, Corea) adoptaron leyes inspiradas en él; lo mismo ocurrió con el reglamento de sustancias químicas REACH (2007)\nocite{Bradford2020Brussels}.
\item \emph{Límites}: en la regulación digital más reciente, algunas empresas han preferido no ofrecer en la Unión ciertos servicios, y Estados Unidos ha respondido con presiones comerciales: la segmentación aparece cuando $F$ es bajo.
\end{itemize}

\paragraph{e) Valoración}

\begin{itemize}
\item \emph{Lo que aporta}: explica con 2 costes un poder que no se ve en las estadísticas del comercio.
\item \emph{Sus límites}:
  \begin{itemize}
  \item el poder depende del tamaño del mercado, y el peso de la Unión en el PIB mundial baja;
  \item normas demasiado estrictas pueden frenar la innovación en la propia Unión, como advirtió el \emph{informe Draghi} (2024) \vertema{3.B.39};
  \item y no sirve donde lo que cuenta es la fuerza: la seguridad, la energía o las materias primas críticas.
  \end{itemize}
\item Por eso la Unión también actúa dentro de las organizaciones internacionales, donde las normas se negocian entre todos.
\end{itemize}

% --------------------------------------------------------------------------
\subsection{Las organizaciones y foros internacionales}\label{sec:organizaciones}

\begin{itemize}
\item \emph{El principio}: la Unión defiende un multilateralismo eficaz basado en reglas, con las Naciones Unidas en el centro (art. 21 TUE), y coopera con las organizaciones internacionales (art. 220 TFUE).
\item \emph{El problema}: casi todas las organizaciones solo admiten Estados. La presencia de la Unión depende de sus competencias (tabla~\ref{tab:organizaciones}):
  \begin{itemize}
  \item \emph{miembro de pleno derecho} donde tiene competencia exclusiva o amplia: la OMC (y los 27 Estados también) \vertema{3.B.33}, la FAO desde 1991 y el G-20, en el que la representan los presidentes del Consejo Europeo y de la Comisión;
  \item \emph{observador} donde la competencia es de los Estados: las Naciones Unidas (observador reforzado desde 2011, con derecho a intervenir en la Asamblea General pero sin voto), el FMI (el BCE es observador en el Directorio) \vertema{3.B.22} y el Banco Mundial, del que es un gran donante \vertema{3.B.30};
  \item \emph{participante sin voto} en la OCDE, a través de la Comisión, y en su Comité de Ayuda al Desarrollo es miembro de pleno derecho \vertema{3.B.19}.
  \end{itemize}
\item \emph{El resultado}: la Unión y sus Estados son los mayores contribuyentes de casi todas estas organizaciones, pero hablan con varias voces. En el FMI, los Estados de la Unión suman en torno a una cuarta parte de los votos, más que Estados Unidos (16~\%), pero se reparten entre varias sillas.
\end{itemize}

\begin{table}[H]
\centering
\caption{La Unión Europea en las organizaciones y foros económicos internacionales}\label{tab:organizaciones}
\small
\begin{tabular}{>{\raggedright\arraybackslash}m{3.6cm}>{\raggedright\arraybackslash}m{3.4cm}>{\raggedright\arraybackslash}m{7.0cm}}
\toprule
\rowcolor{tceeazulpalido} Organización & Estatuto de la UE & Rasgo principal \\
\midrule
Naciones Unidas & Observador reforzado (2011) & Interviene en la Asamblea General; parte en más de 50 convenios \\
OMC & Miembro & La Comisión negocia por los 27, que también son miembros \\
FAO & Miembro (1991) & Vota en lo que es de su competencia \\
G-20 & Miembro & Presidentes del Consejo Europeo y de la Comisión; el BCE en la vía financiera \\
OCDE & Participante sin voto & Miembro del CAD; 22 Estados de la Unión son miembros \\
FMI & Sin estatuto propio & El BCE, observador; los Estados coordinan posiciones (EURIMF) \\
Banco Mundial & Sin estatuto propio & Gran donante de fondos fiduciarios \\
FSB y BPI & Comisión y BCE & Estabilidad financiera y cooperación entre bancos centrales \\
\bottomrule
\end{tabular}
\fuente{elaboración propia a partir del SEAE y de las fichas del Parlamento Europeo\nocite{PE2026FichaOrganizaciones}.}
\end{table}

\begin{anotaciones}
\textbf{Más detalle para las preguntas del tribunal.}
\begin{itemize}
\item \emph{Naciones Unidas}: los Estados de la Unión pagan en torno a una quinta parte del presupuesto ordinario y una parte aún mayor de las contribuciones voluntarias a fondos y programas. Desde la Resolución 65/276 (2011), la Unión puede intervenir en la Asamblea General y presentar propuestas y enmiendas, pero no votar ni presentar candidatos.
\item \emph{FMI}: las decisiones importantes exigen el 85~\% de los votos; Estados Unidos, con algo más del 16~\%, tiene veto, y los Estados de la Unión también lo tendrían si votaran juntos. El EURIMF, creado en 1998, reúne a los directores ejecutivos europeos para coordinar posiciones; el Eurogrupo y el ECOFIN preparan las de la Unión.
\item \emph{G-20}: 19 países, la Unión Europea y, desde 2023, la Unión Africana; reúne en torno al 85~\% del PIB mundial, el 75~\% del comercio y 2 tercios de la población; Alemania, Francia e Italia son miembros, y España es invitado permanente \vertema{3.B.19}.
\end{itemize}
\end{anotaciones}

\begin{itemize}
\item \emph{Recapitulación}: la Unión tiene un marco jurídico propio para su acción exterior, muchos actores y un peso económico enorme; ejerce su poder sobre todo con el mercado, las normas y el dinero. Veamos ahora cómo lo usa con cada grupo de países, empezando por los más cercanos.
\end{itemize}


% ==========================================================================
\section{Visión regional: vecinos, socios y rivales}\label{sec:regional}
\minutos{11}

\begin{itemize}
\item Una vez visto el marco común, pasamos a analizar cómo trata la Unión a cada grupo de países.
\item \emph{Conclusión anticipada}: la Unión organiza sus relaciones en círculos concéntricos (figura~\ref{figu:circulos}). Cuanto más cerca, más integración ofrece y más condiciones pone: a los candidatos, la adhesión; a los vecinos, su mercado y su dinero a cambio de reformas; a los grandes socios, acuerdos y, si hace falta, sanciones. El refrán se cumple: con el vecino cercano la relación es más estrecha y más exigente.
\end{itemize}

\begin{infografia}[H]
\centering
\caption{Los círculos de las relaciones exteriores de la Unión}\label{figu:circulos}
\grafico[0.98\textwidth]{circulos-relaciones}
\fuente{elaboración propia a partir de la Comisión Europea y el SEAE\nocite{CE2026PaqueteAmpliacion,CE2025PactoMediterraneo}. Situación a 14/11/2026.}
\end{infografia}

% --------------------------------------------------------------------------
\subsection{La economía como instrumento geoestratégico}\label{sec:geoestrategico}

\subsubsection{La política de ampliación}\label{sec:ampliacion}

\begin{itemize}
\item \emph{La base}: cualquier Estado europeo que respete los valores del art. 2 TUE puede pedir la adhesión (art. 49 TUE). El Consejo decide por unanimidad, tras consultar a la Comisión y con la aprobación del Parlamento, y el tratado de adhesión lo ratifican todos los Estados (figura~\ref{figu:adhesion}).
\item \emph{Los criterios de Copenhague} (Consejo Europeo de junio de 1993), completados en Madrid (1995):
  \begin{enumerate}
  \item \emph{políticos}: instituciones estables que garanticen la democracia, el Estado de Derecho, los derechos humanos y el respeto de las minorías;
  \item \emph{económicos}: una economía de mercado que funcione y capacidad para hacer frente a la presión competitiva dentro de la Unión;
  \item \emph{el acervo}: capacidad para asumir las obligaciones de la adhesión, con una administración capaz de aplicarlo (Madrid).
  \end{enumerate}
  A ellos se añade la \emph{capacidad de integración} de la propia Unión.
\item \emph{La metodología revisada} (febrero de 2020)\nocite{CE2020Metodologia}: los 33 capítulos que se negocian (35 en total) se agrupan en 6 bloques (\emph{clusters}): fundamentos (Estado de Derecho, economía, instituciones), mercado interior, competitividad, agenda verde, recursos y agricultura, y relaciones exteriores. Los fundamentos se abren los primeros y se cierran los últimos, y el proceso es \emph{reversible}: los retrocesos pueden congelar o reabrir capítulos.
\item \emph{El dinero}:
  \begin{itemize}
  \item el \concepto{IPA III} (Reglamento 2021/1529), 14.162 millones de euros en 2021-2027, ligado a los avances en las reformas\nocite{UE2021IPA};
  \item el \concepto{Plan de Crecimiento para los Balcanes Occidentales}, con su Mecanismo de Reforma y Crecimiento (Reglamento 2024/1449): 6.000 millones (2.000 en donaciones y 4.000 en préstamos) que se pagan cuando el país cumple las reformas de su agenda y que le abren por adelantado partes del mercado único; Moldavia tiene uno propio de 1.900 millones\nocite{CE2026PaqueteAmpliacion};
  \item el \concepto{Mecanismo para Ucrania} (Reglamento 2024/792), 50.000 millones para 2024-2027, del que se habían pagado más de 32.000 millones en octubre de 2026\nocite{CE2026UcraniaDesembolso}, y el \concepto{Préstamo de Apoyo a Ucrania}, 90.000 millones para 2026-2027 (60.000 para defensa y 30.000 de apoyo presupuestario), financiado con deuda de la Unión y acordado por cooperación reforzada en abril de 2026\nocite{CE2026PrestamoUcrania}.
  \end{itemize}
\item \emph{La situación en noviembre de 2026} (paquete de ampliación del 4 de noviembre)\nocite{CE2026PaqueteAmpliacion}:
  \begin{itemize}
  \item 9 candidatos (Albania, Bosnia y Herzegovina, Georgia, Macedonia del Norte, Moldavia, Montenegro, Serbia, Turquía y Ucrania) y un candidato potencial (Kosovo);
  \item \emph{Montenegro} ha cerrado provisionalmente 21 de 33 capítulos y podría cerrar el resto a finales de 2026; ya se redacta su tratado de adhesión, con un paquete financiero de unos 3.200 millones;
  \item \emph{Albania} tiene abiertos los 6 bloques y aspira a terminar en 2027;
  \item \emph{Ucrania y Moldavia} abrieron en junio de 2026 el bloque de fundamentos, cuando Hungría levantó su veto, y han abierto 4 bloques cada una; sus objetivos de cierre son finales y principios de 2028;
  \item la Comisión propone \emph{hojas de ruta} para estos 4 países y un marco de \emph{integración gradual} en energía, transporte, digitalización y defensa;
  \item \emph{Serbia} no avanza, \emph{Bosnia y Herzegovina} se aleja, \emph{Georgia} es candidata «solo de nombre» y con \emph{Turquía} la negociación está parada desde 2018.
  \end{itemize}
\end{itemize}

\begin{infografia}[H]
\centering
\caption{El proceso de adhesión a la Unión Europea}\label{figu:adhesion}
\grafico[0.98\textwidth]{etapas-adhesion}
\fuente{elaboración propia a partir del art. 49 TUE y de la Comisión Europea\nocite{UE2016TUE,CE2020Metodologia}.}
\end{infografia}

\begin{itemize}
\item ¿Por qué hacen los candidatos reformas tan costosas durante tantos años? La respuesta es la condicionalidad.
\end{itemize}

\paragraph{La condicionalidad de la adhesión}\label{sec:condicionalidad}

\subparagraph{a) Idea}

\begin{itemize}
\item La ampliación es la política exterior que mejor ha funcionado: transformó en 15 años las economías y las instituciones de Europa central y oriental. Su mecanismo es la \concepto{condicionalidad}: la Unión ofrece una recompensa muy grande (la adhesión) a cambio de reformas, y solo la da cuando están hechas.
\item \textcite{SchimmelfennigSedelmeier2004} lo formalizaron en el \concepto{modelo de los incentivos externos}: un gobierno adopta las reglas de la Unión si los beneficios de la recompensa superan los costes internos de adoptarlas.
\item Para la economía es, además, un \emph{mecanismo de compromiso}: atar las reformas a un ancla exterior las hace creíbles ante inversores y votantes, como la independencia del banco central ata la política monetaria \vertema{3.A.38}.
\end{itemize}

\subparagraph{b) Modelo}

\begin{itemize}
\item \emph{i) Supuestos}:
  \begin{enumerate}
  \item el gobierno de un país candidato decide si hace una reforma con un coste interno $C$ (rentas que pierden los grupos con poder de veto, coste político);
  \item si la hace, la adhesión llega dentro de $T$ años con probabilidad $p$, la \emph{credibilidad} de la promesa de la Unión; vale $B$ para el gobierno (fondos, mercado, seguridad, votos);
  \item el gobierno descuenta el futuro con un factor $\delta<1$;
  \item la Unión puede dar además recompensas intermedias $g$, ligadas a cada reforma y pagadas antes de la adhesión.
  \end{enumerate}
\item \emph{ii) Desarrollo}:
  \begin{itemize}
  \item \emph{Paso 1: sin recompensas intermedias}, el gobierno reforma si el valor esperado y descontado de la adhesión cubre el coste:
    \funcionclave[eq:condicionalidad]{\text{reforma}\iff C\leq g+p\,\delta^{T}B}
    con $g=0$ en este primer paso (gráfico~\ref{fig:condicionalidad}).
  \item \emph{Paso 2: el plazo}. Si la adhesión se aleja ($T$ sube), $\delta^{T}$ cae y la recta $C=p\,\delta^{T}B$ gira hacia abajo: reformas que antes compensaban dejan de hacerlo. Es la \concepto{fatiga de la ampliación}.
  \item \emph{Paso 3: la credibilidad}. Si la promesa no es creíble ($p$ bajo, por los vetos de un Estado o por las dudas de la Unión), pasa lo mismo.
  \item \emph{Paso 4: las recompensas intermedias}. Con $g>0$, la recta sube en paralelo y el país vuelve a reformar aunque la adhesión siga lejos.
  \end{itemize}
\item \emph{iii) Implicaciones}:
  \begin{itemize}
  \item \uline{la condicionalidad funciona si la recompensa es grande, creíble y cercana, y si el coste interno no es enorme};
  \item explica el éxito de la ampliación de 2004, con una fecha y una promesa creíbles, y el estancamiento de los Balcanes tras 2003 (Salónica), con una promesa lejana y vetos bilaterales;
  \item da sentido a la metodología de 2020 y al Plan de Crecimiento: pagar por reformas antes de la adhesión ($g$) y abrir el mercado único por partes es subir la recta cuando $T$ es grande.
  \end{itemize}
\end{itemize}

\begin{figure}[H]
\centering
\caption{La condicionalidad de la adhesión: cuándo reforma un candidato}\label{fig:condicionalidad}
\grafico[0.6\textwidth]{condicionalidad}
\fuente{elaboración propia a partir de \textcite{SchimmelfennigSedelmeier2004}. El candidato $X$ reforma con la adhesión cercana ($T_1$), deja de hacerlo si se aleja ($T_2$) y vuelve a reformar con recompensas intermedias ($g$).}
\end{figure}

\subparagraph{c) Extensiones}

\begin{itemize}
\item \emph{La inconsistencia temporal}: una vez dentro, la recompensa ya se ha cobrado y la Unión pierde la palanca. «Antes de meter, prometo; y una vez metido, olvido lo prometido», dice el refrán, y algo así ha pasado con el Estado de Derecho en algún Estado miembro tras su adhesión. Por eso la Unión ha añadido herramientas para después: el mecanismo de condicionalidad del Estado de Derecho del presupuesto (2020) y las salvaguardias que propone para los futuros tratados de adhesión.
\item \emph{La geopolítica}: desde 2022 la ampliación se justifica también por la seguridad, lo que sube $B$ para los candidatos y la credibilidad $p$, pero puede tentar a la Unión a rebajar sus exigencias.
\end{itemize}

\subparagraph{d) Evidencia empírica}

\begin{itemize}
\item \emph{Las reformas}: los países de Europa central y oriental adoptaron el acervo más deprisa cuanto más creíble era su adhesión, y los que tenían gobiernos con más que perder (Eslovaquia de Mečiar, Croacia de Tuđman) se retrasaron hasta que cambiaron de gobierno\nocite{Vachudova2005}.
\item \emph{El crecimiento}: con el método del control sintético, \textcite{CamposCoricelliMoretti2019} estiman que la renta por habitante de los países que entraron entre 1973 y 2004 fue de media en torno a un 10~\% más alta en los 10 años siguientes de lo que habría sido sin la adhesión, con la excepción de Grecia.
\end{itemize}

\subparagraph{e) Valoración}

\begin{itemize}
\item \emph{Lo que aporta}: explica con 3 variables (recompensa, credibilidad y plazo) por qué la ampliación transformó a unos países y no a otros.
\item \emph{Sus límites}: trata al candidato como un actor único y racional, sin identidad ni opinión pública; y la credibilidad no es exógena, depende de la unanimidad de los 27.
\item Fuera del círculo de los candidatos, la Unión no puede ofrecer la adhesión: ¿qué condicionalidad le queda con sus vecinos?
\end{itemize}

\subsubsection{La política europea de vecindad}\label{sec:vecindad}

\begin{itemize}
\item \emph{Qué es}: desde 2004, la política con los vecinos que no van a entrar a corto plazo. Su base es el art. 8 TUE: un espacio de prosperidad y buena vecindad con acuerdos específicos.
\item \emph{Su lógica}: «todo menos las instituciones»: acceso al mercado y a los programas de la Unión, y dinero, a cambio de reformas (\emph{más por más}).
\item \emph{2 dimensiones}:
  \begin{itemize}
  \item \emph{la Asociación Oriental} (2009): Armenia, Azerbaiyán, Bielorrusia (suspendida desde 2021), Georgia, Moldavia y Ucrania. Con las 3 últimas hay acuerdos de asociación con zonas de libre comercio de alcance amplio y profundo (ZLCAP, 2014), y hoy son candidatas;
  \item \emph{la Vecindad Meridional}: 10 países del Mediterráneo y Oriente Próximo (Argelia, Egipto, Israel, Jordania, Líbano, Libia, Marruecos, Palestina, Siria y Túnez), con acuerdos de asociación euromediterráneos. Su foro multilateral es la \concepto{Unión por el Mediterráneo} (2008), con 43 miembros y secretaría en Barcelona, heredera del Proceso de Barcelona (1995).
  \end{itemize}
\item \emph{La renovación}:
  \begin{itemize}
  \item en febrero de 2025 la Comisión partió la antigua DG de Vecindad y Ampliación en 2: la \emph{DG ENEST}, para la ampliación y el Este, y la \emph{DG MENA}, para el Mediterráneo, Oriente Próximo y el Golfo, con una comisaria propia para el Mediterráneo\nocite{CE2025DGMENA};
  \item el \concepto{Pacto por el Mediterráneo} (octubre de 2025), con 3 pilares: las personas (educación, movilidad, cultura), unas economías más sostenibles e integradas (energía limpia, conectividad, inversión) y la seguridad, la preparación y la gestión de la migración\nocite{CE2025PactoMediterraneo};
  \item las \emph{asociaciones estratégicas y globales} bilaterales, que juntan comercio, energía, inversión y migración: Túnez (2023), Mauritania (2024, 210 millones), Egipto (2024, 7.400 millones para 2024-2027, sobre todo préstamos macrofinancieros) y Jordania (2025, 3.000 millones)\nocite{Consejo2026Jordania}.
  \end{itemize}
\item \emph{Su valoración}: la condicionalidad es débil (la recompensa es menor que la adhesión y la Unión necesita a sus vecinos para la energía y la migración), y la política se ha ido volviendo más transaccional. España, por su posición, empuja la dimensión mediterránea.
\end{itemize}

\subsubsection{La PESC y las sanciones}\label{sec:sanciones}

\begin{itemize}
\item \emph{La PESC} no es una política económica, pero tiene 2 instrumentos económicos de gran alcance:
  \begin{itemize}
  \item las \concepto{medidas restrictivas}: una decisión PESC por unanimidad (art. 29 TUE) y un reglamento del Consejo por mayoría cualificada (art. 215 TFUE), que aplican las empresas y los Estados; hay más de 40 regímenes en vigor;
  \item el \emph{Fondo Europeo de Apoyo a la Paz} (fuera del presupuesto), que ha financiado material militar para Ucrania.
  \end{itemize}
\item \emph{Las sanciones a Rusia}: desde la invasión de Ucrania (24 de febrero de 2022), 21 paquetes, el último el 23 de julio de 2026\nocite{CE2026Paquete21}:
  \begin{itemize}
  \item \emph{financieras}: inmovilización de unos 210.000 millones de euros del banco central ruso, casi todos en Euroclear, por tiempo indefinido desde diciembre de 2025; exclusión de bancos de SWIFT;
  \item \emph{comerciales}: prohibición de exportar tecnología y bienes de doble uso y de importar acero, carbón, crudo por mar (desde diciembre de 2022) y productos petrolíferos; un mecanismo contra la elusión por terceros países;
  \item \emph{energéticas}: el Reglamento 2026/261 prohíbe el gas ruso: el gas natural licuado (GNL) con contratos a largo plazo desde el 1 de enero de 2027 y el de gasoducto desde el 30 de septiembre de 2027\nocite{UE2026GasRuso};
  \item \emph{contra la flota en la sombra}: más de 550 buques sancionados.
  \end{itemize}
\item La pieza más original es el tope al precio del petróleo, y merece un modelo propio.
\end{itemize}

\paragraph{El tope al precio del petróleo ruso}\label{sec:tope}

\subparagraph{a) Idea}

\begin{itemize}
\item Un embargo total del petróleo ruso habría disparado el precio mundial del crudo y, con él, la inflación en Europa: un golpe para Rusia, pero también para quien sanciona.
\item El G-7 y la Unión buscaron en 2022 un instrumento que quitara a Rusia su renta sin quitar su petróleo del mercado: el \concepto{tope al precio} (\emph{price cap}). Las empresas occidentales, que dominan el transporte marítimo, el seguro y la financiación del crudo, solo pueden dar esos servicios si el crudo ruso se vende por debajo del tope\nocite{JohnsonRachelWolfram2023}.
\end{itemize}

\subparagraph{b) Modelo}

\begin{itemize}
\item \emph{i) Supuestos}:
  \begin{enumerate}
  \item Rusia exporta crudo con un coste marginal bajo $c$ hasta su capacidad $\bar Q$;
  \item el precio mundial $P^*$ lo fija el mercado mundial, y Rusia vende $\bar Q$ porque $P^*>c$;
  \item una parte $\theta$ del crudo ruso necesita servicios occidentales; el resto viaja en una \emph{flota en la sombra} que cuesta $k$ más por barril;
  \item los países que sancionan fijan un tope $\bar P$ con $c<\bar P<P^*$.
  \end{enumerate}
\item \emph{ii) Desarrollo}:
  \begin{itemize}
  \item \emph{Paso 1: sin sanciones}, la renta de Rusia es $(P^*-c)\,\bar Q$ (gráfico~\ref{fig:tope}).
  \item \emph{Paso 2: con tope}, como $\bar P>c$, a Rusia le sigue compensando vender todo $\bar Q$: la oferta mundial no cambia y $P^*$ tampoco.
  \item \emph{Paso 3: el reparto de la renta}. En la parte con servicios occidentales, Rusia cobra $\bar P$; en la otra, $P^*-k$. Su renta es
    \funcionclave[eq:tope]{R=\big[\theta\,\bar P+(1-\theta)(P^*-k)-c\big]\,\bar Q}
    y la diferencia frente a la renta sin sanciones, $\big[\theta\,(P^*-\bar P)+(1-\theta)\,k\big]\bar Q$, pasa a los compradores (en la práctica, la India, China y Turquía, que compran con descuento) o se pierde en la flota en la sombra.
  \end{itemize}
\item \emph{iii) Implicaciones}:
  \begin{itemize}
  \item \uline{el tope quita renta a Rusia sin subir el precio mundial}, a diferencia de un embargo, siempre que $\bar P$ quede por encima de $c$;
  \item su eficacia depende de $\theta$: cuanto más crudo se escapa por la flota en la sombra, menos renta se quita; por eso cada paquete sanciona más buques;
  \item un tope muy bajo es contraproducente: si $\bar P<c$, Rusia dejaría de exportar y el precio mundial subiría.
  \end{itemize}
\end{itemize}

\begin{figure}[H]
\centering
\caption{El tope al precio del petróleo ruso}\label{fig:tope}
\grafico[0.62\textwidth]{tope-petroleo}
\fuente{elaboración propia.}
\end{figure}

\subparagraph{c) Extensiones}

\begin{itemize}
\item \emph{El tope dinámico}: el 18.º paquete (julio de 2025) bajó el tope de 60 a 47,6 dólares por barril desde el 3 de septiembre de 2025, y lo ligó al mercado: se fija un 15~\% por debajo del precio medio del crudo Urales y se revisa cada 6 meses; desde febrero de 2026 es de 44,1 dólares, y el 21.º paquete lo congeló en ese nivel hasta julio de 2027\nocite{CE2026Paquete21}.
\item \emph{La eficacia de las sanciones en general}: los estudios clásicos encuentran que las sanciones logran su objetivo político en torno a un tercio de los casos, más cuando el objetivo es modesto, el sancionado es pequeño y depende del que sanciona, y menos con autocracias grandes\nocite{HufbauerEtAl2007,Pape1997}.
\end{itemize}

\subparagraph{d) Evidencia empírica}

\begin{itemize}
\item \emph{El descuento}: el crudo Urales llegó a venderse en marzo de 2023 unos 32 dólares por barril más barato que el Brent (frente a la paridad de enero de 2022). \textcite{KilianRapsonSchipper2025} atribuyen cerca de la mitad a los mayores costes de transporte y el resto, sobre todo, al poder de negociación de la India; el tope en sí habría pesado poco en el precio de exportación, pero obligó a Rusia a pagar una flota en la sombra.
\item \emph{La renta}: los ingresos del petróleo y el gas del presupuesto federal ruso cayeron de 11,6 billones de rublos en 2022 a 8,8 en 2023, se recuperaron en 2024 y volvieron a caer a 8,5 en 2025; su peso en los ingresos federales bajó del 35,8~\% en 2021 al 22,7~\% en 2025\nocite{MinfinRusia2026}.
\item \emph{La lectura}: el petróleo ruso siguió llegando al mercado y el precio mundial no se disparó, como predice el modelo; la renta rusa cayó, pero menos de lo esperado, porque $\theta$ bajó al crecer la flota en la sombra.
\end{itemize}

\subparagraph{e) Valoración}

\begin{itemize}
\item \emph{Lo que aporta}: un modelo sencillo explica por qué un tope puede ser mejor que un embargo y por qué su punto débil es la flota en la sombra.
\item \emph{Sus límites}: supone que la oferta rusa no responde y que los compradores no se coordinan; ignora la represalia y el coste para Europa del fin del gas ruso, que se pagó con la crisis energética de 2022.
\item \emph{Las sanciones, en conjunto}, han debilitado la economía rusa sin detener la guerra: cumplen una función de contención y de señal más que de coerción.
\end{itemize}

% --------------------------------------------------------------------------
\subsection{El movimiento de bienes y factores:\cortetitulo los grandes socios}\label{sec:socios}

\begin{itemize}
\item \emph{La vertiente comercial predomina}: la Unión tiene la mayor red de acuerdos comerciales del mundo; sus contenidos, el Sistema de Preferencias Generalizadas y la defensa comercial se estudian aparte \vertema{3.B.41}. Aquí interesa el conjunto de la relación con cada gran socio.
\item \emph{Estados Unidos}, el primer socio en bienes y servicios y el mayor inversor:
  \begin{itemize}
  \item en 2025 la relación giró: tras la amenaza de aranceles del 30~\%, el acuerdo político de Turnberry (julio) y la declaración conjunta del 21 de agosto fijaron un arancel máximo del 15~\% para la mayoría de los bienes europeos, a cambio de que la Unión suprimiera sus aranceles industriales\nocite{UE2025DeclaracionEEUU};
  \item en febrero de 2026 el Tribunal Supremo estadounidense anuló los aranceles basados en la ley de emergencias económicas, y la Unión aprobó en junio el reglamento que aplica su parte del acuerdo, que vence en marzo de 2028.
  \end{itemize}
\item \emph{China}, «socio, competidor económico y rival sistémico» desde 2019\nocite{CE2019China}:
  \begin{itemize}
  \item es el primer proveedor de la Unión, con un déficit bilateral récord;
  \item la Unión aplica la \emph{reducción de riesgos} (\emph{de-risking}): diversificar en las materias primas críticas, defensa comercial (derechos a los coches eléctricos, 2024) y control de inversiones;
  \item la cumbre de Pekín (julio de 2025), en el 50.º aniversario de las relaciones, solo acordó un mecanismo para las licencias de exportación de tierras raras.
  \end{itemize}
\item \emph{Rusia}: de proveedor de energía a objeto de sanciones (\ref{sec:sanciones}).
\item \emph{El Reino Unido} y la AELC:
  \begin{itemize}
  \item con el Reino Unido, el Acuerdo de Comercio y Cooperación (2021) y el reajuste de la primera cumbre (19 de mayo de 2025): una asociación de seguridad y defensa y la negociación de un acuerdo sanitario y fitosanitario, la vinculación de los mercados de emisiones y la movilidad de los jóvenes, que se quieren cerrar en la segunda cumbre, el 20 de noviembre de 2026;
  \item con Islandia, Liechtenstein y Noruega, el Espacio Económico Europeo, y con Suiza, acuerdos bilaterales renovados en 2025 \vertema{3.B.35}.
  \end{itemize}
\item \emph{América Latina y el Caribe}:
  \begin{itemize}
  \item una Nueva Agenda (2023) y la cumbre UE-CELAC de Santa Marta (noviembre de 2025), con una agenda de inversiones de Global Gateway de 45.000 millones;
  \item acuerdos con México, Chile y Centroamérica, y con Mercosur: firmado en enero de 2026 y aplicado provisionalmente en su parte comercial desde mayo, a la espera del dictamen del Tribunal de Justicia \vertema{3.B.41}.
  \end{itemize}
\item \emph{África}: la cumbre Unión Africana-Unión Europea de Luanda (noviembre de 2025), en el 25.º aniversario de la asociación, con un paquete de inversiones de Global Gateway de 150.000 millones; y el Acuerdo de Samoa (\ref{sec:historia-coop}).
\item \emph{Asia y el Indo-Pacífico}: la estrategia para el Indo-Pacífico (2021); acuerdos con Japón (2019), Corea, Singapur y Vietnam; las negociaciones con la India, concluidas en enero de 2026; asociaciones de seguridad y defensa con Japón (2024) y Canadá (2025).
\item \emph{El Golfo}: primera cumbre con el Consejo de Cooperación del Golfo (octubre de 2024) y la relanzada negociación de un acuerdo comercial; la segunda, prevista en Riad para octubre de 2026, se aplazó.
\end{itemize}

\begin{itemize}
\item \emph{Recapitulación}: la Unión da a cada círculo lo que puede: adhesión a los candidatos, mercado y dinero a los vecinos, acuerdos o sanciones a los grandes socios. El último círculo, el de los países en desarrollo, tiene su propia política.
\end{itemize}


% ==========================================================================
\section{La política de cooperación al desarrollo}\label{sec:cooperacion}
\minutos{7,5}

\begin{itemize}
\item Una vez vistos los círculos más cercanos, pasamos al más amplio: los países en desarrollo.
\item \emph{Conclusión anticipada}: la Unión y sus Estados son el primer donante del mundo, pero la ayuda europea está cambiando de naturaleza: de las donaciones a las asociaciones y la inversión (Global Gateway), de la generosidad al interés propio declarado, y de crecer a recortarse en varios Estados.
\end{itemize}

% --------------------------------------------------------------------------
\subsection{Consideraciones generales}\label{sec:coop-general}

\subsubsection{Conceptos básicos}\label{sec:conceptos}

\begin{itemize}
\item \emph{La cooperación al desarrollo} es el conjunto de acciones públicas y privadas a favor del desarrollo de otros países; incluye las preferencias comerciales, que no transfieren recursos.
\item \emph{La ayuda oficial al desarrollo} (AOD) es la que cumple los 4 requisitos del CAD de la OCDE: la dan organismos públicos; va a los países de la lista del CAD o a organismos multilaterales; su objetivo principal es el desarrollo económico y el bienestar del receptor; y es concesional. Desde los datos de 2018 los préstamos cuentan solo por su \emph{donación equivalente}, si su elemento de donación supera el 45, el 15 o el 10~\% según la renta del receptor\nocite{OCDE2026DefinicionAOD} \vertema{3.B.29}.
\item \emph{Sus tipos}, según lo que se da:
  \begin{enumerate}
  \item \emph{cooperación financiera}, reembolsable (préstamos blandos) o no (donaciones), incluido el \emph{apoyo presupuestario}, que va al presupuesto del receptor;
  \item \emph{asistencia técnica}: conocimientos, formación, hermanamientos entre administraciones;
  \item \emph{cooperación económica}: proyectos de inversión y ayuda al comercio (\emph{aid for trade});
  \item \emph{ayuda humanitaria y de emergencia}, para salvar vidas tras una catástrofe o un conflicto.
  \end{enumerate}
\item \emph{El objetivo del 0,7~\%} de la RNB lo adoptó la Asamblea General de las Naciones Unidas en 1970; la Unión se ha comprometido a alcanzarlo colectivamente en 2030, con un 0,15-0,20~\% para los países menos adelantados\nocite{Consejo2015FfD}.
\end{itemize}

\subsubsection{¿Por qué una política europea de desarrollo?}\label{sec:por-que}

\begin{itemize}
\item \emph{Por qué es tan difícil salir de la pobreza}: falta capital físico y humano (\autor{Solow}; \autor{Mankiw}, \autor{Romer} y \autor{Weil}) \vertema{3.A.43} \vertema{3.A.44}; hay trampas de pobreza que exigen un empujón coordinado \vertema{3.B.28}; y las instituciones extractivas impiden aprovechar los recursos, como sostienen \autor{Daron Acemoglu}, \autor{Simon Johnson} y \autor{James A. Robinson}\footnote{\autor{Daron Acemoglu}, \autor{Simon Johnson} y \autor{James A. Robinson} recibieron el Premio Nobel de Economía en 2024 «por sus estudios sobre cómo se forman las instituciones y cómo afectan a la prosperidad»\nocite{NobelPrize2024}.} \vertema{3.B.28}.
\item \emph{Por qué la Unión}:
  \begin{itemize}
  \item \emph{equidad}: la diferencia de renta entre los países más ricos y los más pobres es éticamente inaceptable, y los Tratados ponen la erradicación de la pobreza como objetivo;
  \item \emph{eficiencia}: la ayuda es un juego de suma positiva; abre mercados y provee bienes públicos mundiales (salud, clima, estabilidad), y la Unión puede hacerlo a una escala y con una coordinación que no tienen los Estados por separado;
  \item \emph{geopolítica}: da influencia, sobre todo cuando otros donantes se retiran y China ofrece su propia financiación;
  \item \emph{historia}: los lazos con las antiguas colonias de los Estados fundadores, que explican el origen de la política (\ref{sec:historia-coop}).
  \end{itemize}
\item Si la Unión quiere que su ayuda reduzca la pobreza, ¿a quién debe dársela?
\end{itemize}

\subsubsection{El reparto eficiente de la ayuda}\label{sec:reparto}

\paragraph{a) Idea}

\begin{itemize}
\item La ayuda al desarrollo hace crecer sobre todo a los países con buenas políticas, según el estudio de \textcite{BurnsideDollar2000}, muy discutido después \vertema{3.B.29}.
\item \textcite{CollierDollar2002} sacaron la consecuencia normativa: si un donante quiere reducir la pobreza lo más posible con un presupuesto dado, debe repartir la ayuda según la pobreza y la calidad de las políticas de cada país, y no según los lazos históricos o políticos. Es la base de la \concepto{selectividad} de la ayuda.
\end{itemize}

\paragraph{b) Modelo}

\subparagraph{i) Supuestos}

\begin{enumerate}
\item Un donante (la Unión) reparte un presupuesto $\bar A$ entre 2 países, $A_1+A_2=\bar A$.
\item La ayuda reduce la pobreza en cada país con rendimientos decrecientes: la reducción marginal de la pobreza $m_i(A_i)$ es decreciente en $A_i$.
\item $m_i$ es más alta donde hay más pobres y donde las políticas son mejores, porque allí la ayuda se convierte en más crecimiento y el crecimiento llega a más pobres.
\end{enumerate}

\subparagraph{ii) Desarrollo}

\begin{itemize}
\item \emph{Paso 1: el problema}. El donante maximiza la reducción total de la pobreza:
  $$\optim{\max}{A_1,A_2}{\Pi_1(A_1)+\Pi_2(A_2)}{A_1+A_2=\bar A}$$
  donde $\Pi_i$ es la reducción de la pobreza en el país $i$ y $m_i=\Pi_i'$.
\item \emph{Paso 2: la condición de primer orden}. El reparto eficiente iguala la reducción marginal de la pobreza en los 2 países:
  \funcionclave[eq:reparto]{m_1(A_1^*)=m_2(A_2^*)}
  (punto $E$ del gráfico~\ref{fig:reparto}).
\item \emph{Paso 3: el coste de repartir mal}. Si el donante reparte por igual, o por lazos históricos, da de más a un país y de menos a otro, y deja de reducir la pobreza en el triángulo entre las 2 curvas.
\end{itemize}

\begin{figure}[H]
\centering
\caption{El reparto eficiente de la ayuda entre 2 países}\label{fig:reparto}
\grafico[0.62\textwidth]{asignacion-ayuda}
\fuente{elaboración propia a partir de \textcite{CollierDollar2002}. El país 1, con más pobreza y mejores políticas, tiene una curva $m_1$ más alta: en el óptimo $E$ recibe más ayuda que con un reparto igual.}
\end{figure}

\subparagraph{iii) Implicaciones}

\begin{itemize}
\item \uline{La ayuda debe ir a donde reduce más pobreza por euro}: a países pobres y con buenas políticas, y retirarse poco a poco de los de renta media.
\item Con los datos de los 90, \textcite{CollierDollar2002} estimaron que la ayuda sacaba de la pobreza a unos 30 millones de personas al año y que, repartida de forma eficiente, sacaría a unos 80 millones: los países con pobreza severa y políticas aceptables tenían el 74~\% de los pobres y recibían el 56~\% de la ayuda.
\item Para la Unión, explica 2 debates: el peso de la vecindad y de los países de renta media en su ayuda, por razones geopolíticas, y la concentración en los países menos adelantados que piden sus compromisos.
\end{itemize}

\paragraph{c) Extensiones}

\begin{itemize}
\item \emph{La coordinación entre donantes}: si cada donante reparte mirando solo su ayuda, el reparto conjunto es ineficiente y el receptor soporta los costes de tratar con muchos donantes pequeños (\emph{fragmentación}). Coordinar a la Comisión y a los 27 en un solo programa por país (la \emph{programación conjunta} y el Equipo Europa) es aplicar el modelo a todos los donantes europeos a la vez\nocite{KnackRahman2007}. \textcite{BigstenTengstam2015} estiman que concentrar a cada donante europeo en menos países y pasar de proyectos a programas ahorraría muchos costes de transacción y reduciría más la pobreza, aunque choca con los intereses políticos de cada Estado.
\end{itemize}

\paragraph{d) Evidencia empírica}

\begin{itemize}
\item \emph{El resultado de \autor{Burnside} y \autor{Dollar} no es robusto}: con más países y años desaparece\nocite{EasterlyLevineRoodman2004}, y el efecto medio de la ayuda sobre el crecimiento es positivo pero pequeño \vertema{3.B.29}.
\item \emph{El reparto real}: la ayuda de los donantes responde tanto a los lazos coloniales y políticos como a la pobreza y las políticas del receptor\nocite{AlesinaDollar2000}; la de la Comisión va en buena parte a la vecindad y a Ucrania, países de renta media.
\end{itemize}

\paragraph{e) Valoración}

\begin{itemize}
\item \emph{Lo que aporta}: una regla clara (igualar el efecto marginal) que da contenido a la palabra «eficacia».
\item \emph{Sus límites}: depende de medir bien la calidad de las políticas, que es subjetiva; castiga a los países frágiles, donde la pobreza es mayor y las políticas peores; e ignora otros objetivos legítimos de la ayuda (humanitarios, de seguridad o de bienes públicos mundiales).
\item \emph{La lección}: para \autor{William Easterly}, más que repartir mejor desde arriba, la ayuda funciona cuando prueba, evalúa y responde ante los beneficiarios\nocite{Easterly2006}; es lo que pide la agenda de la eficacia de la ayuda, que la Unión ha hecho suya.
\end{itemize}

% --------------------------------------------------------------------------
\subsection{La Unión Europea y la ayuda al desarrollo}\label{sec:aod-ue}

\subsubsection{Relevancia: el primer donante del mundo}\label{sec:relevancia-aod}

\begin{itemize}
\item \emph{El Equipo Europa}: en 2025 los 27 Estados dieron unos 87.000 millones de dólares de AOD, el 0,41~\% de su RNB y la mitad de toda la AOD de los países del CAD; con la parte de las instituciones que no se imputa a los Estados, la cifra es algo mayor\nocite{OCDE2026DAC1,OCDE2026AODPreliminar}.
\item \emph{Las instituciones de la Unión} dieron 23.100 millones de euros (26.100 millones de dólares), un 13,8~\% menos que en 2024 en términos reales; con Alemania, Estados Unidos y el Reino Unido, están entre los mayores donantes del mundo.
\item \emph{Por Estados} (gráfico~\ref{fig:aod-estados}): solo Luxemburgo (0,99~\%), Suecia (0,85~\%) y Dinamarca (0,72~\%) superan el 0,7~\%; Alemania, con el 0,56~\%, fue en 2025 el mayor donante del mundo en volumen, por primera vez por delante de Estados Unidos. Los Estados del Este no llegan al 0,2~\%.
\item \emph{Los recortes}: la AOD de los miembros europeos del CAD bajó del 0,47 al 0,42~\% de su RNB en 2025, con Alemania (del 0,68 al 0,56~\%), Bélgica, Francia, Irlanda y los Países Bajos a la cabeza; Bélgica, Finlandia, Francia, Alemania, los Países Bajos y Suecia han anunciado más recortes para 2025-2027, para pagar más defensa\nocite{OCDE2025Recortes}. España sube (del 0,25 al 0,27~\%). El objetivo colectivo del 0,7~\% en 2030 se aleja.
\end{itemize}

\begin{figure}[H]
\centering
\caption{La AOD de los Estados de la Unión en porcentaje de la RNB, 2025}\label{fig:aod-estados}
\grafico[0.9\textwidth]{aod-ue-estados}
\fuente{elaboración propia a partir de la OCDE\nocite{OCDE2026DAC1}. Donación equivalente; 2025, datos preliminares.}
\end{figure}

\subsubsection{Marco jurídico y objetivos}\label{sec:marco-coop}

\begin{itemize}
\item \emph{La competencia} es paralela (art. 4.4 TFUE): la Unión tiene su política y los Estados la suya, que se complementan y refuerzan mutuamente (art. 208 TFUE).
\item \emph{El objetivo principal} es la reducción y, a largo plazo, la erradicación de la pobreza (art. 208). La Unión debe tenerlo en cuenta en todas sus políticas que afecten a los países en desarrollo: es la \concepto{coherencia de las políticas para el desarrollo} (la PAC o la política pesquera no deberían dañar a los países pobres).
\item \emph{La toma de decisiones}: el Parlamento y el Consejo adoptan los instrumentos por el procedimiento legislativo ordinario (art. 209); la Unión y los Estados coordinan sus programas (art. 210); la ayuda humanitaria tiene su artículo propio (art. 214) y sus principios: humanidad, neutralidad, imparcialidad e independencia.
\item \emph{El Consenso Europeo sobre Desarrollo} (2017), firmado por las 3 instituciones y los Estados, adapta la política a la Agenda 2030 y sus 17 Objetivos de Desarrollo Sostenible, en torno a 5 ejes: personas, planeta, prosperidad, paz y asociaciones\nocite{UE2017Consenso}.
\end{itemize}

\subsubsection{Historia: de Yaundé a Samoa}\label{sec:historia-coop}

\begin{itemize}
\item \emph{La asociación con las excolonias} (figura~\ref{figu:cronologia}):
  \begin{itemize}
  \item el Tratado de Roma (1957) asoció a los países y territorios de ultramar, y en 1959 nació el \emph{Fondo Europeo de Desarrollo} (FED), fuera del presupuesto;
  \item los \emph{Convenios de Yaundé} (1963 y 1969) con 18 Estados africanos independizados;
  \item los \emph{Convenios de Lomé} (1975-2000) con los países de África, el Caribe y el Pacífico (ACP): preferencias comerciales no recíprocas, el STABEX (estabilización de los ingresos de exportación) y el FED.
  \end{itemize}
\item \emph{Una política propia}:
  \begin{itemize}
  \item Maastricht (1992) le dio base jurídica; en 1971 llegó el Sistema de Preferencias Generalizadas para todos los países en desarrollo, y en 2001 la iniciativa \emph{Todo menos armas} para los menos adelantados \vertema{3.B.41};
  \item el \emph{Acuerdo de Cotonú} (2000-2023) sustituyó las preferencias no recíprocas, incompatibles con la OMC, por \emph{acuerdos de asociación económica} recíprocos y asimétricos, y añadió el diálogo político.
  \end{itemize}
\item \emph{El Acuerdo de Samoa} entre la Unión y la Organización de Estados de África, el Caribe y el Pacífico (OEACP, 79 miembros): firmado el 15 de noviembre de 2023 en Apia y aplicado provisionalmente desde el 1 de enero de 2024. Tiene una base común con 6 prioridades (derechos humanos y gobernanza, paz y seguridad, desarrollo humano, crecimiento sostenible, clima y migración) y 3 protocolos regionales (África, Caribe y Pacífico). En noviembre de 2026 todavía no ha entrado en vigor: necesita la ratificación de todos los Estados de la Unión y de 2 tercios de los de la OEACP\nocite{PE2026Samoa}.
\end{itemize}

\begin{infografia}[H]
\centering
\caption{La cooperación al desarrollo de la Unión en el tiempo}\label{figu:cronologia}
\grafico[0.98\textwidth]{cronologia-cooperacion}
\fuente{elaboración propia a partir de la Comisión Europea y del Parlamento Europeo\nocite{PE2026FichaDesarrollo}.}
\end{infografia}

\subsubsection{Los instrumentos}\label{sec:instrumentos-coop}

\begin{itemize}
\item \emph{El IVDCI-Europa Global} (Reglamento 2021/947) reúne desde 2021 casi toda la acción exterior, incluido el antiguo FED: 79.462 millones de euros en 2021-2027\nocite{UE2021IVDCI}:
  \begin{itemize}
  \item \emph{programas geográficos}, 60.388 millones: al menos 29.181 para el África subsahariana y 19.323 para la vecindad, 8.489 para Asia y el Pacífico y 3.395 para América y el Caribe (gráfico~\ref{fig:ivdci});
  \item \emph{programas temáticos}, 6.358 millones (derechos humanos y democracia, sociedad civil, paz y estabilidad, retos mundiales); \emph{respuesta rápida} a crisis, 3.182 millones; y un \emph{colchón} para retos emergentes, 9.534 millones, que se ha usado para los refugiados sirios, Ucrania o la migración;
  \item \emph{objetivos de gasto}: al menos el 93~\% debe ser AOD; el 30~\% para el clima, el 20~\% de la AOD para la inclusión social y el desarrollo humano, y un 10~\% indicativo para la migración. En la vecindad, un 10~\% se reserva a premiar las reformas («más por más»).
  \end{itemize}
\end{itemize}

\begin{figure}[H]
\centering
\caption{El IVDCI-Europa Global 2021-2027 por partidas (millones de euros)}\label{fig:ivdci}
\grafico[0.85\textwidth]{ivdci}
\fuente{elaboración propia a partir del art. 6 del Reglamento (UE) 2021/947\nocite{UE2021IVDCI}. En morado, los programas geográficos (los de la vecindad y del África subsahariana son mínimos).}
\end{figure}

\begin{itemize}
\item \emph{El FEDS+} (Fondo Europeo de Desarrollo Sostenible Plus), dentro del IVDCI: con la Garantía de Acción Exterior, de hasta 53.449 millones, el presupuesto cubre parte del riesgo de los préstamos e inversiones de los bancos de desarrollo, provisionándolo entre el 9 y el 50~\%; se espera que movilice hasta 232.000 millones de inversión en 2021-2027\nocite{CE2026GlobalGateway}. Es el paso de las donaciones a la movilización de inversión privada, con el BEI Global y los bancos de desarrollo de los Estados (en España, COFIDES y la AECID) \vertema{4.A.24}.
\item \emph{El Equipo Europa} (2020, nacido con la respuesta a la covid-19): la Comisión, los Estados, sus agencias y sus bancos de desarrollo (BEI, BERD) presentan una oferta conjunta en cada país, con \emph{iniciativas Equipo Europa} sobre los mismos objetivos. Es la vieja aspiración de coordinar 28 políticas, con más visibilidad.
\item \emph{Global Gateway} (diciembre de 2021), la estrategia de inversión en infraestructuras digitales, de energía y clima, transporte, salud y educación, con 300.000 millones de euros movilizados en 2021-2027 por el Equipo Europa (donaciones, garantías del FEDS+ y préstamos de los bancos europeos):
  \begin{itemize}
  \item en octubre de 2025 la Comisión dio por cumplido el objetivo antes de tiempo, con 306.000 millones movilizados en 2021-2024, y espera superar los 400.000 millones en 2027\nocite{CE2025GlobalGatewayForo};
  \item tiene paquetes regionales, como los de África (150.000 millones) y América Latina y el Caribe (45.000 millones), y proyectos como corredores de transporte, cables submarinos o el hidrógeno verde;
  \item \emph{su crítica}: son cifras de inversión «movilizada» (comprometida y apalancada), no de dinero desembolsado ni adicional, y mezclan ayuda con intereses europeos (materias primas críticas, contratos para empresas europeas).
  \end{itemize}
\item \emph{La ayuda humanitaria}, gestionada por la DG ECHO con organizaciones socias (ONG, agencias de las Naciones Unidas, Cruz Roja): 11.569 millones en 2021-2027 y un presupuesto inicial de 1.900 millones en 2026, que suele crecer durante el año (2.620 millones en 2022)\nocite{ECHO2026}. Se rige por el Reglamento 1257/96 y por el Consenso Europeo sobre la Ayuda Humanitaria (2007).
\item \emph{El comercio}: el Sistema de Preferencias Generalizadas (estándar, SPG+ y Todo menos armas) y los acuerdos de asociación económica \vertema{3.B.41}.
\end{itemize}

\subsubsection{Eficacia, retos y futuro}\label{sec:retos-coop}

\begin{itemize}
\item \emph{La eficacia}: la Unión suscribe los principios de París (2005): apropiación, alineamiento, armonización, gestión por resultados y mutua responsabilidad \vertema{3.B.29}. Sus problemas conocidos son la fragmentación (28 políticas), la lentitud y el uso de la ayuda con fines migratorios.
\item \emph{De la ayuda a las asociaciones}: Global Gateway y el FEDS+ usan la ayuda para atraer inversión, en competencia con la Iniciativa de la Franja y la Ruta china. Es más eficaz para infraestructuras en países de renta media que para los más pobres, donde no hay proyectos rentables que apalancar.
\item \emph{La propuesta para 2028-2034} (julio de 2025): un solo instrumento, \emph{Europa Global}, de unos 200.000 millones de euros (el 0,14~\% de la RNB de la Unión), que fusiona el IVDCI, la preadhesión y la ayuda humanitaria (hasta 25.000 millones), con Ucrania aparte, en una reserva propia\nocite{CE2025EuropaGlobal}:
  \begin{itemize}
  \item se organiza por regiones (Europa ampliada y Este; Oriente Próximo, norte de África y Golfo; África subsahariana; Asia y el Pacífico; América y el Caribe), con un pilar mundial y un colchón mayor;
  \item baja el mínimo de AOD del 93 al 90~\% y quita los suelos de gasto en clima, desarrollo humano y migración: más flexibilidad y más peso de los intereses de la Unión, menos foco en la pobreza\nocite{Bruegel2026EuropaGlobal};
  \item se negocia con el resto del marco, que el Consejo Europeo quiere cerrar a finales de 2026 \vertema{3.B.37}.
  \end{itemize}
\item \emph{España}: dio 4.554 millones de euros de AOD en 2025, el 0,27~\% de su RNB, la cifra más alta en 15 años\nocite{Exteriores2026AOD}. La Ley 1/2023 de Cooperación para el Desarrollo Sostenible y la Solidaridad Global fija el 0,7~\% para 2030 y al menos el 10~\% para la acción humanitaria, y crea el Fondo Español de Desarrollo Sostenible (FEDES), adscrito a la AECID, que sustituye al FONPRODE para la cooperación financiera\nocite{BOE2023LeyCooperacion}; el Plan Director 2024-2027 la desarrolla. La AECID es socia habitual de la Comisión en la cooperación delegada y en las iniciativas Equipo Europa \vertema{4.A.26}.
\end{itemize}

\begin{itemize}
\item \emph{Recapitulación}: la Unión es el primer donante del mundo y tiene una política propia, con un gran instrumento único desde 2021; pero su ayuda se mezcla cada vez más con sus intereses (migración, materias primas, geopolítica), y su volumen depende de unos Estados que recortan.
\end{itemize}

% ==========================================================================
\section*{Conclusión}
\minutos{1,5}

\begin{itemize}
\item \textbf{\emph{Recapitulación}}:
  \begin{itemize}
  \item En conclusión, las relaciones económicas exteriores de la Unión son un conjunto de políticas con una base jurídica común, muchos actores y un poder sobre todo económico: el mercado, las normas, el dinero y la moneda.
  \item Las aplica por círculos: la adhesión condicionada con los candidatos, que funciona si es creíble y cercana; el mercado y el dinero a cambio de reformas con los vecinos; acuerdos y sanciones, como el tope al petróleo ruso, con los grandes socios.
  \item Con los países en desarrollo es el primer donante del mundo, con una política que pasa de la ayuda a la inversión y que debería repartirse donde más pobreza reduce.
  \end{itemize}
\item \textbf{\emph{Relevancia}}:
  \begin{itemize}
  \item En un mundo de bloques y aranceles, la acción exterior es la que decide si la Unión es un actor o un escenario; para España, es la vía de su política hacia el Mediterráneo, América Latina y África.
  \end{itemize}
\item \textbf{\emph{Extensiones}}:
  \begin{itemize}
  \item La política comercial \vertema{3.B.41}; la financiación del desarrollo y la eficacia de la ayuda \vertema{3.B.29}; los bancos de desarrollo \vertema{3.B.30}; el presupuesto y el marco 2028-2034 \vertema{3.B.37}; el papel internacional del euro \vertema{3.B.43}; la cooperación española y sus instrumentos financieros \vertema{4.A.26} \vertema{4.A.24}.
  \end{itemize}
\item \textbf{\emph{Opinión}}:
  \begin{itemize}
  \item La Unión tiene más poder del que usa y menos del que cree: su fuerza son su mercado y sus normas, pero no tiene una sola voz y depende de la unanimidad. Su mejor política exterior ha sido la ampliación, y la apuesta de esta década, Ucrania, Moldavia y los Balcanes, solo saldrá bien si la promesa es creíble y la Unión se reforma para recibirlos.
  \end{itemize}
\item \textbf{\emph{Idea final}}:
  \begin{itemize}
  \item Volvamos al refrán: más vale vecino cercano que hermano lejano. Cuando el hermano lejano sube sus aranceles y recorta su ayuda, la Unión descubre que su fuerza está en sus vecinos y en sus socios.
  \item Por eso su peso económico le sirve para algo más que comerciar: para convertir a sus vecinos en miembros, a sus socios en aliados y a los países pobres en mercados que crecen. Pero solo si cumple lo que promete.
  \end{itemize}
\end{itemize}

\clearpage
\appendix

% ==========================================================================
\section{Anexo: el papel internacional del euro}\label{anexo:euro}

\begin{itemize}
\item \emph{Qué es una moneda internacional}: la que usan fuera de su zona los agentes privados (para facturar, pagar y ahorrar) y los gobiernos (como ancla de su tipo de cambio, para intervenir y como reserva). Son las 3 funciones del dinero \vertema{3.A.35} llevadas al exterior.
\item \emph{Sus determinantes}, en los que la Unión compite con Estados Unidos:
  \begin{enumerate}
  \item el peso de la zona emisora en la producción y el comercio mundiales;
  \item la confianza en la estabilidad de su valor (baja inflación, banco central independiente);
  \item las externalidades de red: cuantos más la usan, más útil es, lo que da mucha inercia a la moneda dominante;
  \item la profundidad y liquidez de sus mercados financieros y la existencia de un activo seguro común.
  \end{enumerate}
  El euro compite bien en los 2 primeros; el dólar le gana en los 2 últimos. Completar la unión bancaria y la de mercados de capitales y emitir más deuda común ayudarían \vertema{3.B.44}.
\item \emph{Ventajas para el emisor}: señoreaje internacional (billetes y deuda que el resto del mundo acepta a cambio de bienes reales), menos riesgo cambiario y costes de transacción para sus empresas, deuda más barata y peso geopolítico (la capacidad de sancionar de Estados Unidos descansa en el dólar). \emph{Inconvenientes}: más volatilidad de la demanda de la moneda (y del tipo de cambio) y la responsabilidad de tener en cuenta el efecto de su política monetaria en el resto del mundo.
\item \emph{Los indicadores en 2025} (BCE, junio de 2026)\nocite{BCE2026PapelEuro}:
  \begin{itemize}
  \item un índice compuesto de uso internacional en torno al 20~\%, unos 5 puntos más que el peso de la zona del euro en el PIB mundial;
  \item el 20,2~\% de las reservas de divisas (dólar, 57~\%), mientras el oro subió al 27~\% de las reservas oficiales totales;
  \item el 24,6~\% de la deuda internacional, el 28,5~\% de las operaciones del mercado de divisas (dólar, casi el 90~\%) y en torno al 23~\% de los pagos por SWIFT;
  \item el 60,0~\% de las exportaciones de la zona del euro al resto del mundo y el 53,1~\% de sus importaciones se facturan en euros.
  \end{itemize}
\item \emph{El euro como ancla} de otros países, de más a menos rigidez (tabla~\ref{tab:anclas}). La zona del euro tiene 21 miembros desde que Bulgaria la adoptó el 1 de enero de 2026.
\end{itemize}

\begin{table}[H]
\centering
\caption{Regímenes cambiarios ligados al euro fuera de la zona del euro, 2026}\label{tab:anclas}
\small
\begin{tabular}{>{\raggedright\arraybackslash}m{4.6cm}>{\raggedright\arraybackslash}m{3.4cm}>{\raggedright\arraybackslash}m{6.0cm}}
\toprule
\rowcolor{tceeazulpalido} Régimen & Estados de la UE & Fuera de la UE \\
\midrule
Euroización (sin moneda propia) & --- & Con acuerdo monetario: Andorra, Mónaco, San Marino y Vaticano; unilateral: Montenegro y Kosovo \\
Caja de conversión & --- & Bosnia y Herzegovina \\
Tipo fijo con bandas (MTC II) & Dinamarca & --- \\
Tipo fijo o estabilizado frente al euro & Rumanía (estabilizado) & Macedonia del Norte; francos CFA (14 países africanos) y CFP; Cabo Verde, Comoras, Santo Tomé y Príncipe \\
Flotación con el euro como referencia & Chequia, Hungría, Polonia & Serbia, Albania; cestas con el euro (Marruecos) \\
Flotación libre con objetivo de inflación & Suecia & --- \\
\bottomrule
\end{tabular}
\fuente{elaboración propia a partir del anexo estadístico del BCE\nocite{BCE2026PapelEuro}.}
\end{table}

\begin{itemize}
\item \emph{El tipo de cambio frente al dólar}: el euro nació en 1999 a 1,17 dólares, cayó a 0,83 en 2000 (burbuja tecnológica en Estados Unidos), subió hasta casi 1,60 en 2008 (déficit exterior estadounidense), osciló en torno a 1,35 en la Gran Recesión, bajó hacia 1,05 con la expansión cuantitativa del BCE (2015-2016) y volvió a la paridad en 2022, cuando la Reserva Federal subió antes los tipos y la guerra de Ucrania golpeó a Europa. El euro flota libremente: un acuerdo cambiario formal con otras monedas lo tendría que decidir el Consejo por unanimidad (art. 219 TFUE).
\end{itemize}

% ==========================================================================
\section{Anexo: el reparto de competencias}\label{anexo:competencias}

\begin{table}[H]
\centering
\caption{Las competencias de la Unión y de los Estados miembros}\label{tab:competencias}
\small
\begin{tabular}{>{\raggedright\arraybackslash}m{4.4cm}>{\raggedright\arraybackslash}m{5.2cm}>{\raggedright\arraybackslash}m{4.4cm}}
\toprule
\rowcolor{tceeazulpalido} Exclusivas (art. 3 TFUE) & Compartidas (art. 4 TFUE) & De apoyo (art. 6 TFUE) \\
\midrule
Unión aduanera; normas de competencia del mercado interior; política monetaria del euro; conservación de los recursos biológicos marinos; política comercial común; celebración de acuerdos internacionales en ciertos casos &
Mercado interior; política social (aspectos del TFUE); cohesión; agricultura y pesca; medio ambiente; consumidores; transportes; redes transeuropeas; energía; libertad, seguridad y justicia; salud pública (aspectos del TFUE); \emph{paralelas}: investigación y espacio, y \textbf{cooperación al desarrollo y ayuda humanitaria} &
Salud humana; industria; cultura; turismo; educación, formación, juventud y deporte; protección civil; cooperación administrativa \\
\bottomrule
\end{tabular}
\fuente{elaboración propia a partir del TFUE\nocite{UE2016TFUE}. La PESC tiene un régimen propio (arts. 23-46 TUE).}
\end{table}

% ==========================================================================
\section{Anexo: de los Objetivos del Milenio a la Agenda 2030}\label{anexo:odm}

\begin{itemize}
\item \emph{Los Objetivos de Desarrollo del Milenio} (2000-2015) concentraron la ayuda en la lucha contra la pobreza con 8 objetivos:
  \begin{enumerate}
  \item erradicar la pobreza extrema y el hambre;
  \item lograr la enseñanza primaria universal;
  \item promover la igualdad de género y la autonomía de la mujer;
  \item reducir la mortalidad infantil;
  \item mejorar la salud materna;
  \item combatir el VIH/sida, el paludismo y otras enfermedades;
  \item garantizar la sostenibilidad del medio ambiente;
  \item fomentar una alianza mundial para el desarrollo.
  \end{enumerate}
\item \emph{La respuesta europea}: en Monterrey (2002) la comunidad internacional se comprometió a dar más ayuda, y la Unión aprobó el \emph{Compromiso de Barcelona} (marzo de 2002): un 0,39~\% de la RNB colectiva en 2006, que se elevó en 2005 al 0,56~\% en 2010 y al 0,7~\% en 2015. La crisis impidió cumplirlo, y en 2015 el plazo pasó a 2030\nocite{Consejo2015FfD}.
\item \emph{La Agenda 2030} (2015) y sus 17 Objetivos de Desarrollo Sostenible son universales, valen también para los países ricos, y unen la erradicación de la pobreza con la sostenibilidad social y ambiental; el Consenso Europeo de 2017 los traslada a la política de la Unión. 2015 fue el Año Europeo del Desarrollo.
\end{itemize}

% ==========================================================================
\section{Anexo: preguntas de test}\label{anexo:test}

\begin{anotaciones}
Preguntas tipo test del primer ejercicio: \href{https://www.victorgutierrezmarcos.es/oposicion/temario/primer-ejercicio/test/simulador.html?tema=3.B.42}{simulador de test}.
\end{anotaciones}

\bibliografia
\end{document}
